% Packages, notation and environments of the qdjj_solver physics manual.
\usepackage[T1]{fontenc}
\usepackage[utf8]{inputenc}
\usepackage{lmodern}
\usepackage[english]{babel}
\usepackage{microtype}
\usepackage{scrhack}
\usepackage{amsmath,amssymb,mathtools,bm}
\usepackage{graphicx}
\usepackage{booktabs,array,tabularx,longtable}
\usepackage{enumitem}
\usepackage[dvipsnames,table]{xcolor}
\usepackage{csquotes}
\usepackage[font=small,labelfont=bf]{caption}
\usepackage{listings}
\usepackage[most]{tcolorbox}
\usepackage[backend=biber,style=phys,biblabel=brackets,maxnames=4,minnames=3,
            doi=true,eprint=true,url=false,isbn=false,pageranges=false]{biblatex}
\usepackage[colorlinks=true,linkcolor=NavyBlue,citecolor=OliveGreen,urlcolor=NavyBlue]{hyperref}
\usepackage[nameinlink,capitalise,noabbrev]{cleveref}
\hypersetup{pdftitle={qdjj\textunderscore solver: physics manual},
            pdfauthor={qdjj\textunderscore solver developers}}

\addbibresource{references.bib}
\graphicspath{{figures/}}
\setlist{itemsep=2pt,topsep=4pt}
\setcounter{tocdepth}{1}
\setlength{\emergencystretch}{1.5em}
\setkomafont{disposition}{\normalfont\bfseries}
\setkomafont{chapter}{\normalfont\bfseries\LARGE}

% --- notation -----------------------------------------------------------------
\newcommand{\code}[1]{\texttt{\detokenize{#1}}}
\newcommand{\pkg}{\textsf{qdjj\_solver}}
\newcommand{\ii}{\mathrm{i}}
\newcommand{\ee}{\mathrm{e}}
\newcommand{\dd}{\mathrm{d}}
\newcommand{\ket}[1]{\lvert #1\rangle}
\newcommand{\bra}[1]{\langle #1\rvert}
\newcommand{\braket}[2]{\langle #1\vert #2\rangle}
\newcommand{\matel}[3]{\langle #1\vert #2\vert #3\rangle}
\newcommand{\expval}[1]{\langle #1\rangle}
\newcommand{\abs}[1]{\lvert #1\rvert}
\newcommand{\norm}[1]{\lVert #1\rVert}
\newcommand{\dg}{^\dagger}
\newcommand{\Tr}{\operatorname{Tr}}
\newcommand{\sgn}{\operatorname{sgn}}
\newcommand{\Heff}{H_{\mathrm{eff}}}
\newcommand{\Hd}{H_{d}}
\newcommand{\Sz}{S^z}
\newcommand{\up}{{\uparrow}}
\newcommand{\dn}{{\downarrow}}
\newcommand{\eps}{\epsilon}
\newcommand{\qe}{\varepsilon}
\newcommand{\Om}{\Omega}

% --- boxes ----------------------------------------------------------------------
\tcbset{boxsep=1.5mm,left=1.5mm,right=1.5mm,top=1mm,bottom=1mm,arc=1mm,boxrule=.5pt,
        fonttitle=\bfseries\small,breakable}
\newtcolorbox{physics}[1][Physical picture]{colback=NavyBlue!4,colframe=NavyBlue!55,title=#1}
\newtcolorbox{incode}[1][In the code]{colback=OliveGreen!4,colframe=OliveGreen!60,title=#1}
\newtcolorbox{caution}[1][Caution]{colback=BrickRed!4,colframe=BrickRed!60,title=#1}

% --- code listings ----------------------------------------------------------------
\lstdefinestyle{python}{language=Python,basicstyle=\ttfamily\small,keywordstyle=\color{NavyBlue}\bfseries,
  commentstyle=\color{gray}\itshape,stringstyle=\color{BrickRed},showstringspaces=false,
  columns=fullflexible,keepspaces=true,frame=single,rulecolor=\color{gray!40},
  backgroundcolor=\color{gray!4},xleftmargin=2mm,xrightmargin=2mm,aboveskip=6pt,belowskip=6pt,
  breaklines=true,literate={−}{{-}}1}
\lstset{style=python}

% Figures are rendered from data/ by `make figures`; show a placeholder if one is missing.
\newcommand{\manualfigure}[1]{\IfFileExists{figures/#1.pdf}{\includegraphics[width=\textwidth]{#1}}%
  {\fbox{\parbox{.9\textwidth}{\centering\vspace{2em}Figure \texttt{#1} not rendered: run \texttt{make figures}.\vspace{2em}}}}}

% --- generated content (tools/gen_docs.py, make quoted) ----------------------------
% Version, pointers to the user manual and the tutorials, and numbers quoted from the figure
% data are generated, so that the text cannot drift from the code, the docs or the data.
% An unknown key is an error rather than a silent "??".
\makeatletter
\newcommand{\qdjj@define}[3]{\@ifundefined{qdjj@#1@#2}{\@namedef{qdjj@#1@#2}{#3}}%
  {\PackageError{qdjj}{#1 key `#2' defined twice}{}}}
\newcommand{\qdjj@use}[3]{\@ifundefined{qdjj@#1@#2}%
  {\PackageError{qdjj}{Undefined #1 key `#2'}{#3}}{\@nameuse{qdjj@#1@#2}}}
\newcommand{\defq}[2]{\qdjj@define{q}{#1}{#2}}
\newcommand{\defusermanual}[2]{\qdjj@define{um}{#1}{#2}}
\newcommand{\deftutorial}[2]{\qdjj@define{tut}{#1}{#2}}
% \q{figure.key}: a number from docs/manual/data, formatted by the figure script's quoted().
\DeclareRobustCommand{\q}[1]{\qdjj@use{q}{#1}{Add it to quoted() in the figure script and run make quoted}}
% \usermanual{docs/page.md}: the page title and path in the user manual.
\newcommand{\usermanual}[1]{User manual: \emph{\qdjj@use{um}{#1}{Run python tools/gen_docs.py}}
  (\texttt{\detokenize{#1}})}
% \tutorial{NN_topic}: the runnable tutorial that accompanies a worked example.
\newcommand{\tutorial}[1]{\par\medskip{\raggedright\noindent\textbf{Tutorial:}
  \emph{\qdjj@use{tut}{#1}{Add the tutorial to examples/catalog.toml and run python tools/gen_docs.py}}
  (\texttt{\detokenize{examples/#1}}).\par}}
\makeatother
% Generated by tools/gen_docs.py from pyproject.toml; do not edit.
\newcommand{\pkgversion}{0.3.0}

% Generated by tools/gen_docs.py from the user-manual headings and examples/catalog.toml.
% Used by \usermanual{path} and \tutorial{directory}; do not edit.
\defusermanual{README.md}{qdjj\_solver}
\defusermanual{CONTRIBUTING.md}{Contributing}
\defusermanual{docs/bath_representations.md}{Bath and reservoir representations: accuracy, convergence, and cost}
\defusermanual{docs/chain_expansion.md}{Padé chain-expansion reservoirs}
\defusermanual{docs/custom_models.md}{General impurities and custom operators}
\defusermanual{docs/developer.md}{Architecture and extension guide}
\defusermanual{docs/dmrg.md}{DMRG states and bond-dimension refinement}
\defusermanual{docs/documentation.md}{Documentation policy and tooling}
\defusermanual{docs/floquet.md}{Floquet calculations}
\defusermanual{docs/getting_started.md}{Installation and a first checked calculation}
\defusermanual{docs/index.md}{qdjj\_solver documentation}
\defusermanual{docs/interfaces.md}{Using the solver and saving results}
\defusermanual{docs/models.md}{Physical models and conventions}
\defusermanual{docs/numerics.md}{Numerical methods and convergence}
\defusermanual{docs/qp_convergence.md}{Convergence with the quasiparticle cutoff}
\defusermanual{docs/qp_workflow.md}{QP calculations and observable-specific convergence}
\defusermanual{docs/references.md}{Citation and method references}
\defusermanual{docs/releases.md}{Release artifacts and procedure}
\defusermanual{docs/reservoir_workflow.md}{Choosing and refining superconducting reservoirs}
\defusermanual{docs/spectroscopy.md}{Probe matrix elements and spectra}
\defusermanual{docs/static_calculations.md}{Static states, observables and transitions}
\defusermanual{docs/troubleshooting.md}{Troubleshooting}
\defusermanual{docs/user_manual.md}{PDF manuals: reading, building and publication}
\defusermanual{docs/validation.md}{Validation and development}
\defusermanual{docs/workflows.md}{Parameter scans, saved results and reproducible analysis}
\defusermanual{docs/reference/api.md}{API reference}
\defusermanual{docs/reference/cli.md}{Command-line reference}
\defusermanual{docs/reference/concordance.md}{From the equations to the code}
\defusermanual{docs/reference/configuration.md}{Calculation input format}
\defusermanual{docs/reference/conventions.md}{Conventions}
\defusermanual{docs/reference/models.md}{Model and drive construction}
\defusermanual{docs/reference/observables.md}{Observables and probes}
\defusermanual{docs/reference/options.md}{Option reference}
\defusermanual{docs/reference/records.md}{Saved records and artifacts}
\defusermanual{docs/reference/results.md}{Result reference}
\deftutorial{00_first_calculation}{A first calculation}
\deftutorial{01_subgap_states}{Subgap states, singlet and doublet}
\deftutorial{02_driven_dot}{The gate-driven isolated dot}
\deftutorial{03_two_level_system}{The proximitized dot as a driven two-level system}
\deftutorial{04_high_frequency}{The high-frequency limit}
\deftutorial{05_convergence}{Numerical convergence of a driven calculation}
\deftutorial{06_phase_drive}{The phase-driven junction}
\deftutorial{07_gate_drive}{The gate-driven junction}
\deftutorial{08_two_tone}{Two-tone spectroscopy of the pair transition}
\deftutorial{09_photon_assisted_tunnelling}{Photon-assisted tunnelling into a driven junction}
\deftutorial{10_spin_resonance}{Electron spin resonance of the doublet}
\deftutorial{11_dmrg_convergence}{DMRG bond-dimension convergence}

\InputIfFileExists{generated/quoted}{}{\PackageError{qdjj}{generated/quoted.tex missing}{Run make quoted}}
